\section{おわりに}
本報告書では，災害や地理による通信障害に影響されない通信の開発とその可視性の確保を目的に，ArduinoとLED，スピーカ，振動モータを用いた無線通信システムを開発した．
具体的には，PCの文字入力から照度センサでLEDの光を検知し，スピーカから出力された音をマイクでその周波数を検出し，加速度センサで振動モータの振動を判断して文字を伝送した．
3つの通信を結合したシステムでの通信は可能であったが，通信精度や通信速度の面では実用には程遠かった．
加えて，どの通信も実施可能性や精度が環境に左右されやすく，自分たちで最適な環境を用意する必要があったことからも，本システムは有効であるとはいえないということがわかった．
通信可能な環境を増やしたり，センサの性能を最大限に引き出せれば，有効なシステムになる可能性がある．
現在通信できている環境での周囲の光や雑音の影響の調査をすることで，これらの問題の解決につながるだろう．